\thispagestyle{empty}

\begin{singlespace}

{\small
  \begin{center}
    \begin{minipage}{0.8\textwidth}
      \begin{center}
        {\normalsize ОТЗЫВ}\\[1em]
        на дипломный проект\\студента факультета информационных технологий 
        и управления\\Учреждения образования <<Белорусский государственный университет информатики и радиоэлектроники>>\\
        Былькова Даниила Владимировича\\
        на тему: Интеллектуальная система мониторинга рисков и обеспечения безопасности путешественников
      \end{center}
    \end{minipage}
  \end{center}

На время дипломного проектирования перед студентом Быльковым~Д.\,В. была поставлена задача разработать интеллектуальную систему мониторинга рисков и обеспечения безопасности путешественников.
Тема является весьма актуальной, так как в условиях современной глобальной геополитической и климатической нестабильности классические реактивные подходы уступают место автоматизированным средствам проактивного мониторинга и экстренного оповещения.
Проблема фильтрации шума в неструктурированных новостных потоках и точного пространственного сопоставления зон рисков с маршрутами пользователей является ключевой при создании подобных платформ.

На основе комплексного анализа профильной литературы и международных стандартов Быльков~Д.\,В. определил оптимальные подходы к интеллектуальному анализу неструктурированных текстовых данных.

В рамках проектирования разработаны событийно-ориентированная микросервисная архитектура системы. Программный комплекс реализован на базе современных объектно-ориентированных и скриптовых языков программирования с привлечением специализированных фреймворков декларативной разработки, а в качестве хранилища данных задействована объектно-реляционная СУБД с расширением для обработки пространственных объектов.

Разработанный конвейер анализа событий и созданное программное обеспечение являются результатом комплексного исследования предметной области. Работа демонстрирует умение систематизировать научно-техническую литературу и успешно применять фундаментальные теоретические знания для решения актуальных прикладных задач. Студент Быльков~Д.\,В. проявил себя как самостоятельный исследователь, способный проводить глубокий системный анализ и реализовывать сложные наукоемкие ИТ-проекты.

Работа над проектом велась ритмично и в строгом соответствии с утвержденным календарным графиком.
Пояснительная записка и графический материал оформлены аккуратно, грамотно и в полном соответствии с требованиями нормативных документов.

Система, разработанная в рамках дипломного проекта, внедряется в отделы по туристической и корпоративной безопасности нескольких крупных организаций.

Дипломный проект Былькова~Д.\,В. полностью соответствует техническому заданию, отличается глубокой научной и инженерной проработкой темы и выполнен с применением современных ИТ-технологий.

Считаю, что Быльков~Д.\,В. в полной мере освоил процессы и технологии проектирования интеллектуальных систем, подготовлен к самостоятельной работе по специальности 1-40~03~01
<<Искусственный интеллект>> и заслуживает получения квалификации <<Инженер-системотехник>>.

  \vfill
  \noindent
  \begin{minipage}{0.54\textwidth}
    \begin{flushleft}
      Руководитель дипломного проекта:\\
      инженер-программист\\
      
      \underline{\hspace*{2em}} \underline{\hspace*{6.5em}} \the\year{}~г.
    \end{flushleft}
  \end{minipage}
  \begin{minipage}{0.44\textwidth}
    \begin{flushright}
      \underline{\hspace*{3cm}} Д.\,С.~Кисляк
    \end{flushright}
  \end{minipage}
}

\end{singlespace}

\clearpage